Use the Laplace transform to reduce the partial differential equation
\begin{equation}
\frac{\partial^2 u}{\partial r^2}
+
\cos(\varphi)\frac{\partial u}{\partial \varphi}
=
r^2
\label{50000033:pde}
\end{equation}
on the plane described in polar coordinates $(\varphi,r)$
with boundary conditions
\begin{equation}
\left.
\begin{aligned}
u(0,r)                                   &= -1     &\qquad &\forall r>0    \\
\frac{\partial u}{\partial \varphi}(0,r) &= e^{-r} &\qquad &\forall r>0    \\
u(\varphi,0)                             &= 3      &\qquad &0<\varphi<2\pi \\
\frac{\partial u}{\partial r}(\varphi,0) &= 0      &\qquad &0<\varphi<2\pi
\end{aligned}
\quad\right\}
\label{50000033:boundary} \end{equation}
to a family of ordinary differential equations with suitable boundary
conditions.
Do the boundary conditions~\eqref{50000033:boundary} uniquely determine
the solutions of the equation~\eqref{50000033:pde}?

\begin{hinweis}
Remember that in polar coordinates, the function $u(\varphi,r)$ must
satisfy periodic boundary conditions: $u(\varphi+2\pi,r)=u(\varphi,r)$.
\end{hinweis}

\begin{loesung}
The transform has to be done with respect to the $r$-coordinate,
we write the transformed coordinate as $s$.
The transform of the differential equation becomes
\begin{align*}
-
\frac{\partial u}{\partial r}(\varphi,0)
-
su(\varphi,0)
+
s^2\mathscr{L}u(\varphi,s)
+
\cos(\varphi)
\frac{\partial}{\partial \varphi} \mathscr{L}u(\varphi,s)
&=
\frac{2}{s^3}
.
\intertext{Writing $\mathscr{L}u(s,\varphi)=U_s(\varphi)$,
this equation becomes}
-
3s
+
s^2 U_s(\varphi)
+
\cos(\varphi)U'_s(\varphi)
&=
\frac{2}{s^3}
\\
\cos(\varphi)U'_s(\varphi)
+
s^2U_s(\varphi)
&=
3s
+\frac{2}{s^3},
\end{align*}
which is a family of linear first order ordinary differential
equations with parameter $s$.
It also satisfies the boundary conditions obtained from transforming
the first two boundary conditions in \eqref{50000033:boundary}
\begin{align*}
\mathscr{L}u(0,s)
&=
-\frac1s
&&\Rightarrow&
U_s(0)&=-\frac{1}{s}
\\
\frac{\partial \mathscr{L}u}{\partial y}(0,s)
&=
\frac{1}{s+1}
&&\Rightarrow&
U'_s(0)
&=
\frac{1}{1+s}.
\end{align*}
As the boundary conditions uniquely determine the function $U_s(y)$,
there can only be a single solution for $u(x,y)$.
\end{loesung}

\begin{diskussion}
The Laplace transform of the function $r^2$ may not be as well known
as the other transforms, here is its derivation using partial
integration twice:
\begin{align*}
(\mathscr{L}r^2)(s)
&=
\int_0^{\infty} r^2e^{-sr}\,dr
=
\biggl[
-\frac{r^2e^{-sr}}{s}
\biggr]_0^\infty
+
\int_0^\infty
\frac{2re^{-sr}}{s}
\,dr
=
-\frac{2}{s}
\int_0^\infty
re^{-sr}\,dr
\\
&=
-\frac{2}{s}
\biggl(
\biggl[
-
\frac{re^{-sr}}{s}
\biggr]_0^\infty
-
\int_0^\infty
-
\frac{e^{-sr}}{s}
\,dr
\biggr)
=
\frac{2}{s^2}
\int_0^\infty e^{-sr}\,dr
=
\frac{2}{s^2}
\biggl[ -\frac{e^{-sx}}{s}\biggr]_0^\infty
=
\frac{2}{s^3}.
\end{align*}
\end{diskussion}

\begin{bewertung}
Choice of coordinate for the Laplace transform ({\bf C}) 1 point,
Laplace transform of derivatives ({\bf D}) 1 point,
substitution of boundary conditions ({\bf S}) 1 point,
Laplace transform of the equation ({\bf E}) 1 point,
Laplace transform of boundary conditions ({\bf B}) 1 point,
uniqueness of solution ({\bf U}) 1 point.
\end{bewertung}

